\subsection{Accounting for the eight historical reductions}

The historical sequence mixes changes to the register program with changes
to its binary implementation.  The following accounting is therefore by
mechanism and surviving role, rather than by numerical state names alone.
The original reductions and the changed K230 frontend are attributed to
their respective sources.\cite{proveit235,k230}

\begingroup
\small
\renewcommand{\arraystretch}{1.12}
\begin{longtable}{@{}p{0.16\linewidth}p{0.23\linewidth}p{0.53\linewidth}@{}}
\toprule
Historical step & Relevant K correspondence & Disposition in this construction\\
\midrule
\endfirsthead
\toprule
Historical step & Relevant K correspondence & Disposition in this construction\\
\midrule
\endhead
\bottomrule
\endfoot
$243\to242$; old $27\to15$
& Old $15\mapsto K11$, $32\mapsto K26$
& Already in K230: $K11/0=0LK26$ and the corresponding caller redirects.
No additional saving.\\
$242\to241$; old $39\to17$
& Old $17\mapsto K13$, $106\mapsto K91$, $37\mapsto K31$
& Already in K230: the two $K13$ rows move left into $K91$ and $K31$.
No additional saving.\\
$241\to240$; old $90\to72$
& Old $90\mapsto K79$, $72\mapsto K63$, $142\mapsto K121$
& Transplanted after the total transfer-loop rewrite. Delete $K79$ into
$K63$; four exceptional positive blocks have lengths $9,3,3,3$.\\
$240\to239$; delete old $218$
& Old $218\mapsto K175$, $179\mapsto K150$; projection role becomes $K65$
& Transplanted after source sharing and an additional K-specific entry-tag
repair. Split the two entry contexts of $K175$ between $K150$ and $K65$.\\
$239\to238$; old $169\to80$
& Old $169\mapsto K140$, $80\mapsto K71$, $91\mapsto K80$
& Transplanted with coordinated helper reuse. One local block needs an
extra right-neighbor guard; the final machine is checked directly.\\
$238\to237$; old $235\to98$
& Old $98\mapsto K85$, $240\mapsto K186$
& No direct counterpart for old 235. Its mapped children would be
$(K180,K12)$; the nearby K199 instead has $(K180,K166)$.\\
$237\to236$; old $223\to182$
& Projection role $K65$, retained $K152$, helpers $K159,K169$
& Transplanted after two source rotations repair the helper layout and
the retained state's entry tag. Delete $K65$ into $K152$.\\
$236\to235$; old $193\to74$
& Required mapped pairs $(K157,K172)$ and $(K99,K182)$
& Neither pair is a current K decoder node. The earlier correspondence
old $74\mapsto K65$ predates their different repurposing; it gives no
further fold of the final table.\\
\end{longtable}
\endgroup

Thus two historical savings were already included, four are newly
transplanted, and two have no direct surviving pair.  The accepted order is
\[
 230\xrightarrow{\text{delete }K79}229
 \xrightarrow{\text{delete }K140}228
 \xrightarrow{\text{delete }K175}227
 \xrightarrow{\text{delete }K65}226.
\]
Source reassociations and repairs between these arrows add no working
state.  The complete correspondence is retained in
\code{comparison/TRANSPLANT_CLASSIFICATION.md}.

\subsection{The first two folds and their distinct proof scopes}

The total macro replacement proved in Lemma~\ref{lem:macro} changes the
two decoder successors of $K121$ from $(K31,K63)$ to $(K63,K13)$.
It applies at four source occurrences.  Then delete $K79$, redirect its
incoming transitions to $K63$, and set
\[
 K63/0=0LK123,\qquad K63/1=1LK30.
\]
The lost state distinction is recovered from already present predecessor
bits.  At old $K63$ the two bits immediately left are $10$; at old $K79$
the immediately preceding bit is $1$.  These facts are inductive from
the common initial state, without restricting its initial tape.
The four exceptional blocks reproduce the old write, displacement,
successor and entire tape in $9,3,3,3$ transitions.  The other 456 old
rows commute with the state identification in one transition.  Both the
generic symbolic checker and a separate four-case checker prove this fold.

For the second fold, identify $K140$ with $K71$ and change
\[
 K71/0=0LK91,\qquad K80/1=1LK138.
\]
The changed helper $K80$ is part of the construction, so this is more
than replacing two identical rows.  The exceptional old rows
$K71/0,K140/0,K140/1,K80/1$ have matching new blocks of lengths
$5,3,5,9$.  Their predecessor guards follow from the binary table,
except that $K140/1$ additionally requires a $1$ immediately to the right.
The local certificate explicitly records this condition.  It does not
claim an all-tape theorem for that fold.  The complete runtime proof
below establishes the final machine's behavior on every valid frame,
including this helper's actual use, independently of extrapolating
the conditional local argument.

\subsection{A source schedule that leaves usable decoder contexts}

The first source reassociation changes the fragment
\[
 K85=(K19,K65),\quad K65=(K114,K12)
 \;\longrightarrow\;
 K85=(K98,K12),\quad K98=(K19,K114).
\]
Here a displayed pair describes source children.  The paired subtree
$K98$ already exists, so the new left child requires no new decoder.
The freed label $K65$ is reused as a backward projection helper, with
$K175/0=0LK65$ and helper rows $0RK12,1RK86$.
Sharing a source node does not by itself justify these backward entries;
the binary runtime is a separate obligation.

The next reassociation is also K-specific.  Originally
$K202=((K149,K150),K151)$, with inner state $K201$.
The right rotation makes $K202=(K149,(K150,K151))$.
Thus $K201$ reaches $K150$ on its zero branch, repairing the entry tag
needed below.  The exact children are 149,150,151: common ancestry with
an older source is not a license to substitute a superficially similar
subtree.

Write $L_{q,r}$ for $(A,(B,C))\mapsto((A,B),C)$ at outer state $q$
and inner state $r$, and $R_{q,r}$ for the inverse.  The complete ordered
schedule after the loop replacement is
\[
\begin{gathered}
 L_{85,65},\ R_{202,201},\ L_{174,119},\ L_{148,154},\
 R_{52,67},\ R_{67,229},\\
 R_{210,174},\ R_{174,101},\ R_{226,224},\ L_{205,134},\
 R_{207,205},\ L_{101,209},\ L_{101,208}.
\end{gathered}
\]
The certificate identifies every occurrence path, so shared fragments
are rewritten consistently.  The first four stages temporarily give a
maximum zero run of seven.  The final schedule restores six, retains
5,121 leaf occurrences and 211 distinct source DAG nodes, and has maximum
path depth 41.  This matters because primitive operations modify path
bits before reset: a smaller source DAG alone does not preserve the
binary frame invariant.

Stages 5 and 6 cross the terminal AX break.  In their right rotation
$((A,B),C)\mapsto(A,(B,C))$, the third chunk is
$C=\mathrm{CB}(1)$ and is \emph{not closed}: it has one escaping break.
A rule allowing only three closed chunks would exclude both moves.
Their admissibility comes from comparing the complete source control
graph, including every break target and both decrement branches.
These moves permit the eventual zero-run repair without another reset
state, but they also move a fused machine instruction and therefore
require the callback repair described below.

\subsection{The context split and the final three-step fold}

In the prepared 228-state basis, all entries to $K150$ write a zero
and move right.  The two callers $K155/1,K176/1$ enter $K175$ after
writing one; $K184/0$ enters it after writing zero.  Moreover, old $K65$
is entered only inside the backward $K175/0$ projection.
Delete $K175$ by representing its left-zero context with $K65$ and its
left-one context with $K150$.  Old $K150$ retains its left-zero context.
The literal changes are
\[
\begin{aligned}
 K150/1&=1LK173,& K155/1&=1RK150,& K176/1&=1RK150,\\
 K65/0&=0RK19,& K65/1&=1RK61,& K184/0&=0RK65.
\end{aligned}
\]
Six exact open-tape identities handle the two reads in these three
contexts.  Old blocks have length one or three; new blocks have length
one, three or five.  Each has the same whole tape and head at its
endpoint.  Ordinary rows commute, and the sole entry to old $K65$
ensures it is an internal macro state rather than an omitted checkpoint.
This proves halting equivalence from state 0 on every initial binary tape.

The rotations $L_{174,119}$ and $L_{148,154}$ supply the final fold's
entry tag and helper layout.  Old $K65$ now has left bit zero, whereas
$K152$ has left bit one.  Delete $K65$ into $K152$ and set
\[
 K152/0=0LK159,\qquad K152/1=1LK169.
\]
The four exceptional paths are
\[
\begin{array}{ccl}
 \text{old state/read}&\text{left bit}&\text{new state path}\\
 K65/0&0&K152,K159,K166,K19\\
 K65/1&0&K152,K169,K99,K61\\
 K152/0&1&K152,K159,K118,K155\\
 K152/1&1&K152,K169,K148,K171.
\end{array}
\]
All are three-transition blocks with unchanged surrounding tape.
The checker proves all 454 old rows: 450 one-step blocks and these four
three-step blocks.  No extra state is charged for either helper tree.

\subsection{Repairing the relocated break without adding a state}

The AX rotations move the terminal break path from
\texttt{111010010011} to \texttt{11101001001111}.
The old fused action $K52/1=0LK27$, followed by $K27/1=0LK40$,
clears two trailing ones.  Copying it to $K229/1$ does not clear
the four ones in the new path.  There is a second issue: the old
backward callback $K103/1=0LK52$ would enter the newly repurposed
decoder instead of its former primitive action.  The concrete failures
occurred at different CB1 leaves, so repairing just one was insufficient.

An existing duplicate, $K18/1=0LK27$, resolves both problems:
\begin{verbatim}
103 1 0 L 18
229 1 0 L 103
\end{verbatim}
The original two-step callback $K103,K52,K27$ becomes
$K103,K18,K27$ with exactly the same tape, head and endpoint.
The complete clearing chains are now
\[
 K103\to K18\to K27\to K40\to K4,
 \qquad
 K229\to K103\to K18\to K27\to K40\to K4.
\]
An inductive invariant of the complete repaired table supplies the
left-to-right predecessor patterns $011$ at $K103$ and $0111$ at
$K229$.  On reading one, the first chain clears three trailing ones
and their preceding zero; the second clears four and their preceding
zero.  Both enter the ordinary prefix-increment routine $K4$ on the
preceding cell.  Every unspecified tape cell remains free in these
local proofs.  A negative control explicitly rejects the stale
two-one action on the four-one suffix.

The resulting machine, rather than an assumed sum of compatible folds,
is the object of the complete 6656-case runtime proof and exact startup
check.  The source replay, these local identities, and that unbounded
runtime proof form ordinary mathematical certificates available for
internalization in ZFC.  They establish the new transport used in
Theorem~\ref{thm:main}; they are not a claim of a new Rocq kernel check
for the 226-state table.
